% Generated by scripts/generate_lean_blueprint.py; do not edit by hand.
\chapter*{Formalization blueprint}
\addcontentsline{toc}{chapter}{Formalization blueprint}
This blueprint maps every theorem-like statement in the manuscript to its curated Lean counterpart. A linked declaration is not automatically a complete formalization: only entries whose ledger status is exactly \texttt{full} receive the \texttt{leanok} marker. Partial, conditional, special-case, and open interfaces remain visibly incomplete.

\chapter{Part A}

\begin{conjecture}[The target]\label{conj:target}
\lean{Conjecture.CONJ_target}
{{The conjecture, stated precisely, in the normal form the programme
attacks.}}
\end{conjecture}
\begin{remark}
\textbf{Lean alignment.} \texttt{CONJ\_\allowbreak{}target}. Ledger status: \texttt{statement only}. Placeholder: the target conjecture is stated (not proved) in lean/\allowbreak{}conjectures.lean under the CONJ\_\allowbreak{} prefix. Only entries with status exactly 'full' receive \textbackslash{}leanok.
\end{remark}

\begin{theorem}[Example placeholder]\label{thm:example}
\lean{Conjecture.CONJ_target}
{{Replace with the first real reduction. Delete this environment once the
manuscript has content.}}
\end{theorem}
\begin{remark}
\textbf{Lean alignment.} \texttt{CONJ\_\allowbreak{}target}. Ledger status: \texttt{statement only}. Placeholder entry for the example theorem environment in master.tex; replace both together with real content. Every manuscript label must have exactly one entry in this ledger.
\end{remark}
